\section{Methodology and Implementation}
This chapter presents the design principles and implementation of
\textbf{Cedar}, a prototype language and compiler for describing
C layouts declaratively.
The goal of the project is to give programmers
\textbf{greater trust} that in-memory representations of data structures
corresponds to their intent.
Cedar achieves this bya introducing a small external language
for layout declarations, paired with a static checker that validates 
those declarations against C type definitions and generates
layout consistent C code.

While we previously reviewed related work, this chapter focuses on what
was built: the theoretical foundations on which Cedar stands, the kind
of layouts it supports, and the design of its grammar and compilation
pipeline.
The work draws on two main sources of inspiration—\textbf{CompCert} and
\textbf{Dargent}—whose ideas on form the conceptual basis for
Cedar's semantics and tooling.

\subsection{Foundations: CompCert and Dargent}
Cedar's design is built upon the semantic foundations provided by
\textbf{CompCert}~\cite{leroy2009formal}, and the declarative layout
principles introduced by \textbf{Dargent}~\cite{dargent}.
Understanding these two systems explains both the motivation
and structure of Cedar.

\subsubsection{CompCert verified C semantics as a foundation}
CompCert~\autocite{compcert} is a formally verified compiler of the C language
Its significance lies 
in \textbf{formal definition of the C memory model}~\autocite{leroy2012compcert},
which specified how C objects are represented as contiguous blocks of bytes,
how pointer arithmetic and aliasing work, and under what condition two
memory states can be considered equivalent.
Although Cedar is not yet built \textit{within} CompCert, it adopts
the same conceptual views on C memory: objects are typed, laid out in
addressable bytes, and subjected to well-defined alignment and
padding rules.

By grounding Cedar's reasoning in this model, layout checking can be 
interpreted consistently within verified C semantics.
For example, when Cedar computes bit-ranges, and checks for non-overlapping
fields, it effectively enforces that the resulting layout 
would be \textbf{well-formed under CompCerts memory abstraction}. The
connection ensures that, if Cedar layouts were eventually compiled through
CompCert, their structural assumptions would remain sound.

\subsubsection{Dargent: declarative layouts and type-layout separation}
While CompCert provides the semantic bedrock, Cedar's language design
follows the philosophy of \textbf{Dargent}~\autocite{dargent}.
Dargent introduced a declarative approach to specifying how algebraic data types
are represented in memory, separating the physical layout from the logical
structure.

Cedar adopts key ideas—\textbf{type-layout separation} and \textbf{explicit
field placement}—but targets them for plain C rather than Cogent. Unlike Dargent,
Cedar does not require a restricted functional language, or a theorem prover;
it provides lightweight static checking that enforces the same
well-formedness properties. These adaptations transform Dargent's research
concept into a more practical, C-oriented tool.

\subsubsection{Bridging the Two influences}
Cedar can thus be viewed as a \textbf{bridge between CompCert and Dargent}:
From CompCert, it inherits a precise understanding of how C objects
occupy memory; from Dargent, it inherits a declarative notation for describing
that occupation explicitly.
Cedar's goal is to close the gap between declarative layout specifications
in self-contained languages, and the C ecosystem, producing verifiable code
that remain compatible with existing C compilers.

The remainder of this chapter builds on these foundations.
We continue by describing the subset of C types supported by Cedar and how
layouts are represented. Later we detail the language patterns and features.
Finally, we summarize the compiler architecture that realizes it.

%Figure of Cedar's lineage
\subsection{Language Overview}
Cedar operates over a well-defined subset of the C type system that
reflect the representations defined by CompCert. The
prototype focuses on data structures whose layout is
statically determinable: \textbf{structures, arrays, and
tagged unions}. These cases cover the majority of layout-sensitive
code in low-level programming while remaining tractable for static
checking.

\subsubsection{Structures}
Cedar treats C \texttt{struct} as ordered collections of named
fields occupying contiguous regions of memory.
Each field has a known size and alignment derived from its
base type.
A layout declaration for a structure describes the physical
order and spacing of its fields using placement
directives
such as \texttt{at}, \texttt{before}, and \texttt{after}.
During checking, Cedar ensures that:
\begin{itemize}
    \item All fields of the corresponding C \texttt{struct}  are
    represented exactly once,
    \item No two fields overlap in their bit ranges, and
    \item The overall size and alignment conform to the
    base types requirements.
\end{itemize}

This model allows Cedar to detect and properly handle changes that could
otherwise be silently mishandled by compiler-specific padding or mis-ordered
field declarations—one of the most common causes of data-layout bugs in C.

\subsubsection{Arrays} 
Arrays are represented as repeated layouts of a uniform elemenet type.
Cedar supports both fixed-length and statically known arrays.
The layout checked verifies that the element layout
is well-formed and then replicates its size and alignment
across the declared array length.
Array support enables Cedar to express memory regions such as
device register tables, or network buffers that consist of repeated
records.

\subsubsection{Primitive and Derived Types}
Base types such as \texttt{char}, \texttt{short}, \texttt{int}, \texttt{long},
and their fixed-width equivalents (\texttt{int8\_t}, \texttt{uint32\_t}, etc.)
are treated as atomic elements with known bit sizes and alignment constraints.
Cedar assumes the same type sizes as the host compiler, which can be
configured when the tool is built.
Pointer types are recognized but not yet given explicit layouts;
they are treated as opaque machine words, leaving pointer
validity and aliasing to the underlying C semantics.
Future extensions could integrate CompCert's pointer-value
model to reason about pointer fields more precisely.

\subsubsection{Layout Representation}
Internally, every Cedar layout is reduced to a \textbf{set of
non-overlapping bit-ranges} withing a contiguous address space.
Each range corresponds to a field array element and is annotated with its
logical name, type, and alignment.
FIGURE illustrates this mapping for a simple \textttt{struct} example.

%TODO Figure of example
% layout Student at 0 {}
% name: char[32]
% id: uint43__t after name
% grade: uint8_t after id

The compiler computer explicit offsets—\texttt{name} [0-255],
\texttt{id} [256-287], \texttt{grade}[288-295]—and verifies
that these regions are contiguous and non-overlapping.

This representation follows the same notion of "memory chunk" used in CompCert,
where each object occupies a sequence of bytes and field addresses are
computed as offsets within that chunk.
By maintaining this correspondence, Cedar's checked
layouts can be compiled through any C toolchain, including CompCert,
without violating its memory-safe assumptions.

\subsubsection{Scope and Limitations}
The current prototype omits constructs whose semantics 
complicate static layout reasoning, such as: bitfields,
flexible array members, recursive layouts, and untagged unions.
These features will be considered in future work
once the core checking mechanisms are stable.
Despite these restrictions, the implemented subset is expressive
enough to capture common data representations used in
device drivers, network protocols, and binary file formats.

\subsection{Cedar Language Design and Core Features}
Cedar extends C with a minimal set of declarative constructs for specifying
how data structures are laid out in memory.
Its syntax follows C's lexical convnetions but introduces
new top-level \texttt{layout} blocks
that describe the physical
placement of fields.
These layouts declarations are currently written in separate
\texttt{.cedar} source file and linked to existing C type
definition.
The language is intentionally small: its expressiveness
comes from composable placement directives rather than
complex new forms.

\subsubsection{Core Grammar}
FIGURE summarizes the grammar of Cedar's core language, adapted from
Anderson's proposal~\autocite{cedar_george}.

%figure of the grammar George

Each \texttt{layout} declaration associates a name (matching a C type) with
a collection of fields. Each field declaration specifies the type and
placement of a field relative to the layout's origin
or to other fields.

The \texttt{placement} clause may be absolute (\texttt{\@}) or relative
(\texttt{before}/\texttt{after}).
Recursive or parametric layouts are excluded.

%Example of student layout
%layout Student = record {
%  name  : U8[32] @ 0B,
%  id    : U32    after name,
%  grade : U8     after id
%} @ 0B;

In this example, the \texttt{Student} layout begins at offset 0.
The field \texttt{name} is an array of 32 characters
occupying the first 256 bits; \texttt{id} is placed immediately
after it, and \texttt{grade} follows \texttt{id}.
Offsets are computed automatically during the layout-checking phase.
If two fields were to overlap, or if a field were missing compared
to the associated C type, the checked
would reject the layout.

%Figure annotated same example as above.

\subsubsection{Language Semantics}
Cedar's semantics are deliberately simple and static.
Each layout defines a total mapping from logical fields to \textbf{bit-ranges}
within a contiguous address space.
The compiler computes these ranges deterministically based on field
sizes and placement directives.

\textbf{Key properties enforced by the checked:}
\begin{itemize}
    \item \textbf{Completeness} — every field of the corresponding C type
    must appear in the layout.
    \item \textbf{Non-overlap} — no two fields occupy intersecting bit ranges 
    \item \textbf{Order Consistency} — relative directives (\texttt{before}/\texttt{after})
    must refer to existing fields.
    \item \textbf{Alignment} — offsets must respect the alignment constraints
    of their base types.
\end{itemize}

A layout satisfying these rules is \textit{well-formed}. The checking
algorithm ensures that every accepted layout could correspond to a valid C
memory representation.

\subsubsection{Design Considerations}
Cedar's syntax supports declarative rather than arithmetic.
Instead of writing numeric offsets directly, programmers may
express structural relationships—"\texttt{id} after \texttt{name}"—that
the compiler resolves into precise addresses.
This improves readability and reduces the risk of human error, while
still allowing absolute placement (\texttt{@}) when required
for interoperability with binary formats and hardware registers.

To keep the semantics tractable, Cedar omits features like computed
field expressions or conditional layouts.
The resulting language is small enough to reason about statically and
expressive enough to describe practical C structures.

\subsubsection{Supported Directives Summary}
\begin{center}
    \begin{tabular}{|| c c c||}
        \hline
        Directive & Purpose & Example \\ [0.5ex]
        \hline\hline
        \texttt{@} & Absolute placement.  & \texttt{id : u32 @ 64B;} \\
        \hline
        \texttt{after} & Place field immediately after another. & \texttt{grade : u8 after id;} \\
        \hline
        \texttt{before} & Place field immediately before another. & \texttt{prefix : u8 before name;} \\
        \hline
        \texttt{layout} & Define a new layout block. & \texttt{layout Student { ... }} \\
        \hline
    \end{tabular}
\end{center}

Together , the grammar and these directives define Cedar's core language.
They form the user-facing layer of the system which is then translated by
the compiler into a sequence of well-defined intermediate representations.
These representations capture progressively more information
about the layout, allowing the checked to validate and later generate code
for the specified structures.

\subsection{Compiler Architecture}
Cedar's compiler is organized as a sequence of \textbf{pure translation stages}
that progressively enrich the information available about each layout.
This implementation is lightweight and research-oriented.

%Figure about compiler flow

Each stage transforms the program into a slightly more concrete form,
preserving the structure while adding derived information such as absolute
offsets and alignment constraints.

\subsubsection{Parsing and Resolution}
The grammar is parsed into a structure abstract syntax tree (\texttt{L}).
A resolution pass then substitutes symbolic references (e.g., \texttt{after name})
with explicit dependencies between fields.
Each field returns its declared attributes—type, byte offset, relative placement,
and optional endianness.
At this point the layout is still declarative: field positions are
described relative to one another rather than as numeric offsets.

The parser currently supports: Nested constructs, quoted field names, and forward
references between fields. Comments and diagnostics, and multiple layouts per
field are not yet implemented. Parse errors are reported generically.

\subsubsection{Static Checking}
In the next stage (LR), the compiler resolves \textbf{relative placements} and
\textbf{nesting}. Every symbolic reference (\texttt{after x}, \texttt{before y})
is replaced with a concrete byte offset computed from previously defined fields.
Nested \texttt{record} blocks are \textbf{flattenned} into a single
linear sequence of fields using compound names (e.g., \texttt{grade_maths}).
Endianness and element size information are propagated through this transformation.

The resulting \textit{lowered representation (LR)} is a canonical, order-independent
list of fields, each with a fixed, offset, size and endianness tag.

The \textbf{checked representation (CR)}  introduces explicit byte ranges.
Offsets are multiplied by eight to produce bit ranges for each field, which
allow the code generator to reason about sub-word accessess and multi-word
values.
Cedar performs basic consistency checks during this phase:

\begin{itemize}
    \item fields do not overlap,
    \item offsets must be non-negative and monotonically increasing, and
    \item all fields must have defined sizes.
\end{itemize}

Full alignment checking, total-size verification, and overlap proofs are
not yet implemented, but the structure is designed to accommodate them later
\subsubsection{Code Generation}
The final stage translated the checked representation into compile-able C code.
The output consists of:
\begin{enumerate}
    \item A header file defining the a flat array of \texttt{uint32_t} words;
    and
    \item A C source file containing \textit{accessor} and \textit{mutator}
    functions for each field.
\end{enumerate}

Nested records are emitted as flattened field names; arrays are unrolled into
repeat accessors. Per-field endianness annotations are honoured by inserting
byte-swap operations in generated setters and getters.

No runtime support or special compiler modifications are required. The generated
code can be compiled by any C compiler, including \textbf{GCC} and \textbf{CompCert}.

\subsubsection{Design Rationale}
The architecture emphasizes two main principles:
\begin{enumerate}
    \item \textbf{Transparency}. Each representation (L, LR, CR) can
    be inspected, making the compiler's reasoning traceable
    \item \textbf{Modularity}. Additional analyses or transformations
    can be inserted between stages without disrupting others.
    \item \textbf{Extensibility}. The separation of passes makes it straightforward
    to extend the language (for example, by adding union layouts or precise bit-field control)
    without disrupting the core pipeline.
\end{enumerate}

This design allow Cedar to be extended in the future with
additional features, and potentially support for automated
verification, similar to Dargent.

It provides a practical foundation for Cedar's prototype: small enough
to experiment with, buth structured in a way that preserves the formal
clarity of its inspirations.

\subsection{Achievements and Scope}
The Cedar prototype demonstrated that declarative data-layout specification
can be integrated directly into the C programming environment in a practical
way.
Its development produced a complete end-to-end toolchain—from parsing and static
checking of layout descriptions to the generation of layout-consistent C
code—that validates the design concepts outlined in this chapter.

\subsubsection{Implemented Features}
The current implementation supports:

\begin{itemize}
    \item \textbf{Declarative layout defintions} for \texttt{structs},
    fixed-size arrays, and tagged unions.
    \item \textbf{Placement directives} (\texttt{at}, \texttt{before}, \texttt{after})
    that determine field ordering and spacing.
    \item \textbf{Automatic offset computation} and total-size calculation.
    \item \textbf{C code generation} producing equivalent declarations
    and accessor functions.
\end{itemize}

Each of these features we tested on representative examples derived from
Dargent's examples and from common systems-programming idioms such as
device registers. The prototype successfully detected layout mismatches
and generated compilable C code for all supported constructs.

\subsubsection{Scope and Limitations}
Cedar intentionally restricts its expressiveness to keep checking decidable
and the implementation lightwieght.
The prototype does \textbf{not} yet support:
\begin{itemize}
    \item Bit-fields or overlapping union members;
    \item Flexible array members or dynamically sized layouts;
    \item Nested or recursive layout definitions;
    \item Automated diagnostic recovery;
\end{itemize}

These omissions do not reflect missing functionality but a deliberate research
boundary: the implemented subset is sufficient to evaluate Cedar's design
principles and verify that they operate coherently within C's type system.
The language design and compiler architecture are modular enough that
these features could be added in future versions without
major restructuring.

\subsubsection{Practical Outcome}
The resulting system is a functioning compiler that demonstrates:
\begin{enumerate}
    \item Declarative layout syntax can describe C data structures succinctly;
    \item Static checking can verify structural inconsistencies that ordinary
    C compilers accept silently;
    \item Generated code remains portable and compatible with existing toolchains.
\end{enumerate}

Together, these results validate Cedar's underlying hypothesis: that
separating layout intent from C type declarations can improve
programmers' confidence in the correctness of low-level data representations.

The next chapter presents concrete example and results, illustrating how
Cedar's layout translate into verified C code and how its checker behaves
on representative cases.